\section{Related Work}
\label{leffa:sec:related_work}

\subsection{Controllable Person Image Generation}

\noindent \textbf{Appearance Control}.
Virtual try-on seeks to automate outfit changes in person images without distorting garment details, which is a longstanding challenge~\citep{yang2020towards, choi2021viton, ge2021disentangled, ge2021parser, he2022style}.
Recently, diffusion-based methods have emerged as the de facto solutions~\citep{zhu2023tryondiffusion, kim2024stableviton, xu2024ootdiffusion, wan2024improving, shen2024imagdressing}.
Some pioneering works~\citep{morelli2023ladi, gou2023taming, zeng2024cat}
enhance fine-grained detail preservation by incorporating additional modules, such as a textual inversion module with CLIP~\citep{radford2021learning} for enhanced conditions, a warping model to merge garments with garment-agnostic person images, and ControlNet~\citep{zhang2023adding} with DINOv2~\citep{oquab2023dinov2} to extract garment details.
Further improvements~\citep{kim2024stableviton, choi2024improving, xu2024ootdiffusion, wan2024improving} leverage complex model structure designs, external models and textual information for enhanced performance.
Recently, \citet{chong2024catvton} achieve detail preservation without extra models but requires complex training strategies~\citep{zhou2024dream} to reduce the training instability.

\noindent \textbf{Pose Control}.
Similarly, pose transfer presents a challenge in capturing the intricate structure of the pose transformation while preserving the fine-grained details of the textures~\citep{ma2017pose, siarohin2018deformable, li2019dense, zhu2019progressive, ren2020deep, men2020controllable, zhang2021pise, albahar2021pose, zhou2021cocosnet, sanyal2021learning, zhang2022exploring, zhou2022cross, ren2022neural, shen2023advancing, shen2024imagpose}. 
\citet{bhunia2023person} first leverage diffusion models with a texture diffusion module for quality enhancement.
\citet{han2023controllable} further propose a pose-constrained attention module and replace skeletons with dense pose~\citep{guler2018densepose} for more accurate pose conditioning.
\citet{lu2024coarse} introduce a perception-refined decoder to progressively generate person images from coarse to fine while employing hybrid-granularity attention to enhance detail.
Recently, \citet{pham2024cross} replace UNet~\citep{ronneberger2015u} with the masked diffusion transformer~\citep{peebles2023scalable, gao2023masked}, achieving efficient and high-quality generation.

Although the existing methods have achieved significant progress on both virtual try-on and pose transfer tasks, most of them rely on complex module designs and extra models for fine-grained detail preservation.
In contrast, we introduce a method that ensures stable training without complex designs, reducing detail distortion without adding any inference cost or extra parameters.


\subsection{Attention in Controllable Generation}

The attention mechanism, a crucial module in generative models, has been extensively studied and shown to significantly enhance model generation quality when effectively guided~\citep{liu2024faster, xie2023boxdiff, phung2024grounded, guo2024focus, chefer2023attend, cao2023masactrl, xiao2024fastcomposer, hertzprompt}.
How to improve attention mechanism has also been studied in controllable person image generation.
Some methods~\citep{ren2020deep, zhou2021cocosnet, ren2022neural, tang2020xinggan} propose customized attention mechanisms.
For instance, \citet{ren2020deep} develop a differentiable global-local attention to reassemble inputs at the feature level, while \citet{zhou2021cocosnet} leverage iterative patch matching to enhance attention quality.
\citet{ren2022neural} apply neural texture extraction and distribution, using semantic filters to improve attention maps and image quality.
On the other hand, several methods directly optimize attention maps via supervision.
\citet{han2023controllable} propose a pose-constrained attention loss to model appearance-pose interactions.
\citet{kim2024stableviton} introduce an attention total variation loss to enhance focus on garment regions, though it lacks explicit guidance to attend to correct regions.
In contrast, we propose a regularization loss that guides attention maps to attend to correct regions without additional annotations, enhancing attention quality and reducing fine-grained detail distortion.


\subsection{Developments Since Publication}
\label{leffa:sec:post_publication}
\leffa{} appeared at CVPR 2025. Three developments since then bear directly on the claims of this chapter: a change of backbone, a move away from masks, and a reassessment of how this task should be evaluated.

\noindent \textbf{From U-Nets to diffusion transformers.}
The dual-U-Net architecture that \leffa{} is built on has largely been replaced by diffusion transformer backbones. Voost~\citep{lee2025voost} trains a single DiT jointly on try-on and try-off, so that each garment-person pair supervises both directions, and reports state-of-the-art results without task-specific networks or auxiliary losses. This matters for the present chapter in a specific way. \leffa{} is a loss on an attention map, and we showed it to be model-agnostic across U-Net-based generators. Whether it transfers to a DiT, where attention is global over tokens rather than structured by a convolutional hierarchy, is untested here. The mechanism has no obvious dependence on the backbone, since any architecture in which target tokens attend to reference tokens admits the same flow-field interpretation, but this remains an expectation rather than a result.

Voost also suggests an alternative route to the same goal. \leffa{} improves correspondence by supervising the attention map directly, whereas joint try-on and try-off training improves it by giving each pair two supervisory directions, so that the model must learn a correspondence it can invert. These two approaches are complementary rather than competing. Whether combining explicit attention supervision with bidirectional training compounds the benefit or saturates it is an open and inexpensive experiment.

\noindent \textbf{Mask-free generation.}
Section~\ref{leffa:sec:limitations} identifies the dependence on garment segmentation as a limitation and proposes a mask-free formulation as the natural extension. Others have since carried out that extension. JCo-MVTON~\citep{wang2025jco} removes the human body mask entirely, injects control signals through dedicated conditional pathways in a multi-modal DiT, and handles garment-person correspondence through fusion in the self-attention layers, with refined positional encodings and attention masks for spatial alignment. This resolves the failure mode noted in Sec.~\ref{leffa:sec:limitations}, in which a mis-segmented collar or sleeve produces a visibly wrong composition however well correspondence is resolved. It supersedes part of the pipeline used here rather than the contribution itself, because attention-level correspondence supervision and mask-free conditioning address different components.

\noindent \textbf{Evaluation.}
The concern raised in Sec.~\ref{leffa:sec:limitations}, that FID is a distributional statistic largely blind to the local errors this chapter targets, has since been addressed systematically. VTBench~\citep{hu2025vtbench} decomposes try-on evaluation into five dimensions, including texture preservation, complex-background consistency, and hand-occlusion handling, and grounds them in human preference annotations. Its stated motivation is that current metrics fail to reflect human perception in unpaired settings and that existing test sets are limited to indoor scenes. Both criticisms apply to the evaluation in this chapter. A texture-preservation dimension measures what \leffa{} is designed to improve far more directly than FID does, and the results reported here are likely to understate the effect. Any reassessment of this method should use such a decomposed protocol rather than the aggregate metrics used at the time.

\noindent \textbf{Assessment.}
The central idea, that attention maps in reference-conditioned diffusion can be read as correspondence fields and supervised as such, has not been superseded. The mechanism remains a cheap addition to a training pipeline with no inference cost. What has been superseded is the surrounding apparatus: the U-Net backbone, the segmentation dependency, and the evaluation protocol. A reader adopting this work today should take the loss and replace the pipeline.
